% Part: proof-theory
% Chapter: natural-deduction
% Section: sequents

\documentclass[../../../include/open-logic-section]{subfiles}

\begin{document}

\olfileid{pt}{nat}{seq}

\olsection{সিকোয়েন্টে স্বাভাবিক নিষ্পাদন: \Log{N2c} ও~\Log{N2i}}

\Log{N1c} বা \Log{N1i}-তে !!^a{proof}~$\delta$
দেখায়, তার অন্তিম-!!{formula}~$!A$ খোলা অনুমানগুলি
থেকে অনুসৃত হয়। $\Gamma$ যদি সেই !!{formula}s~$!B$-এর
সমষ্টি হয়, যাদের $!B^x$ প্রমাণ~$\delta$-তে খোলা
অনুমান হিসেবে ঘটে, তাহলে লিখতে পারি
$\delta \Proves \Gamma \Sequent !A$। এ থেকে
বোঝা যায়, সিকোয়েন্টের জন্যও একটি স্বাভাবিক
নিষ্পাদন-পদ্ধতি গড়া যায়। তাতে প্রতিটি সিকোয়েন্ট
ডান দিকের !!{formula}টি কোন খোলা অনুমানগুলির উপর
নির্ভর করে, অর্থাৎ ওই সিকোয়েন্টে শেষ হওয়া
উপ-!!{proof}-এ কোনগুলি খোলা, সেই তথ্য বহন করে।
এটিই সিকোয়েন্টে, বা “সিকোয়েন্ট-রীতিতে”,
স্বাভাবিক নিষ্পাদন; ফলিত পদ্ধতিদুটির নাম
\Log{N2c} ও~\Log{N2i}।
% BN-SRC-555 TARGET-CORRECTION-BEGIN
\par\smallskip\noindent\textnormal{\small\textbf{উৎস-সংশোধন (BN-SRC-555):} এই অনুচ্ছেদের শুরুতে উৎসে G1c/G1i লেখা; আলোচিত সূত্রবৃক্ষ ও পাশের N1/N2 তুলনা অনুযায়ী N1c/N1i করা হয়েছে।}
% BN-SRC-555 TARGET-CORRECTION-END

\Log{N1} ও \Log{N2}-এর মিল আরও স্পষ্ট করতে,
বাঁ দিকের সমষ্টি~$\Gamma$-তে $x:!B$ রূপের
চিহ্নযুক্ত !!{formula}s রাখি। \Log{N1}-এর
বিধিগুলি কেবল পূর্বধারণার !!{formula}s-এ,
অর্থাৎ অব্যবহিত উপ-!!{proof}s-এর অন্তিম
!!{formula}s-এ কাজ করে। তাই \Log{N2}-এর
বিধিগুলিও শুধু সিকোয়েন্টের ডান দিকে কাজ করে।
সিকোয়েন্টের বাঁ দিকের চিহ্নযুক্ত সূত্রগুলিকে
সংক্ষেপে \emph{প্রসঙ্গ} বলতে পারি।

\Log{N2}-এর প্রমাণে অনুমানের বদলে সর্বোচ্চে
থাকে এই রূপের স্বতঃসিদ্ধ-সিকোয়েন্ট:
\[x:!A \Sequent !A.\]
বিধিগুলি \Log{N1}-এর বিধিগুলির সঙ্গে হুবহু
মেলে এবং \olref{tab:N2}-তে দেওয়া আছে।

\subfile{rules-N2}

\Log{N1c} ও \Log{N1i}-এর মতো এখানেও
\Intro{\lif}, \Intro{\lnot}, \Elim{\lor}
ও \Elim{\lexists} বিধিতে অনুমান অবমুক্ত করা
যেতেও পারে, নাও পারে। তাই \Log{N2c}
ও~\Log{N2i}-তে সংশ্লিষ্ট বিধির সঠিক প্রয়োগ
তখনও অনুমোদিত, যখন অব্যবহিত পূর্বধারণার
প্রসঙ্গে সংশ্লিষ্ট !!{formula}s $!A$, $!B$
ও $!A(a)$ অনুপস্থিত।

\begin{prop}
$\delta$ যদি \Log{N1c} বা \Log{N1i}-তে
$!A$-এর !!a{proof} হয়, তাহলে \Log{N2c}
(\Log{N2i})-তে $\Gamma \Sequent !A$-এর
!!a{proof}~$\delta'$ আছে, যেখানে $\Gamma$
হলো এমন সব~$x:!B$-এর সমষ্টি, যাদের জন্য
$!B^x$ হলো~$\delta$-র খোলা অনুমান।
\end{prop}

\begin{proof}
  $n=\pheight{\delta}$-র উপর আরোহ করি।
  $n=0$ হলে $\delta$ কেবল একটি খোলা
  অনুমান~$!A^x$। তখন $\Gamma = \{x:!A\}$
  এবং $\Gamma \Sequent !A$ হলো
  \Log{N2c} বা~\Log{N2i}-এর স্বতঃসিদ্ধ।

  $n>0$ হলে $\delta$-র শেষ অনুমানবিধি অনুযায়ী
  ক্ষেত্র আলাদা করি। এখানে শুধু একটি ক্ষেত্র
  আলোচনা করি; বাকিগুলি অনুশীলন।

  ধরি, শেষ বিধি $\Intro{\lif}$। তাহলে
  $\delta$-র শেষাংশ
  \[
  \AxiomC{$\Discharge{!B}{x}$}
  \RightLabel{$\delta_1$}
  \DeduceC{$!C$}
  \DischargeRule{\Intro{\lif}}{x}
  \UnaryInfC{$!B \lif !C$}
  \DisplayProof
  \]
  আরোহের অনুমানে \Log{N2c} (\Log{N2i})-তে
  $\Gamma' \Sequent !C$-এর !!a{proof}~$\delta_1'$
  আছে, যেখানে $\Gamma'$ হলো~$\delta_1$-এ
  খোলা থাকা চিহ্নযুক্ত অনুমানগুলির সমষ্টি।
  $!B^x$ যদি $\delta_1$-এর খোলা অনুমান হয়,
  তাহলে \Intro{\lif} বিধিতে সেটি অবমুক্ত
  হয়; ফলে~$\delta$-র খোলা অনুমানের সমষ্টি
  $\Gamma = \Gamma' \setminus \{x:!B\}$।
  অন্য কথায়, $\delta_1'$ হলো
  $x:!B, \Gamma \Sequent !C$-এর !!a{proof};
  তাই $\delta'$ নেওয়া যায়
  \[
  \AxiomC{}
  \RightLabel{$\delta_1'$}
  \Deduce$x:!B, \Gamma \fCenter !C$
  \RightLabel{\Intro{\lif}}
  \UnaryInf$\Gamma \fCenter !B \lif !C$
  \DisplayProof
  \]
  % BN-SRC-556 TARGET-CORRECTION-BEGIN
  \par\smallskip\noindent\textnormal{\small\textbf{উৎস-সংশোধন (BN-SRC-556):} এই N2 প্রমাণবৃক্ষের পূর্বধারণায় উৎসে x:B লেখা, অথচ প্রসঙ্গের সূত্রটি সর্বত্র !B; x:!B করা হয়েছে।}
  % BN-SRC-556 TARGET-CORRECTION-END
  $!B^x$ যদি~$\delta_1$-এর খোলা অনুমান না হয়,
  তাহলে \Intro{\lif} কোনো অনুমান অবমুক্ত
  করে না; $\delta$ ও $\delta_1$-এর খোলা
  অনুমান একই। \Log{N2c} ও~\Log{N2i}-তে
  \Intro{\lif}-এর পূর্বধারণায় প্রসঙ্গসূত্র
  $x:!B$ থাকা বাধ্যতামূলক নয়; তাই
  $\delta'$ নেওয়া যায়
  \[
  \AxiomC{}
  \RightLabel{$\delta_1'$}
  \Deduce$\Gamma \fCenter !C$
  \RightLabel{\Intro{\lif}}
  \UnaryInf$\Gamma \fCenter !B \lif !C$
  \DisplayProof
  \]

  $\delta$-র শেষে অন্য বিধি থাকলে
  ক্ষেত্রগুলি অনুরূপ।
\end{proof}

\begin{prop}
$\delta$ যদি \Log{N2c} বা \Log{N2i}-তে
$\Gamma \Sequent !A$-এর !!a{proof} হয়,
তাহলে \Log{N1c} (\Log{N1i})-তে
অন্তিম-!!{formula}~$!A$-এর !!a{proof}~$\delta'$
আছে, যার খোলা অনুমান $!B^x$—প্রতিটি
$x:!B \in \Gamma$-র জন্য।
\end{prop}

\begin{proof}
  আবার~$\delta$-র !!{height}~$n$-এর উপর
  আরোহ করি। $n=0$ হলে $\delta$ একটি
  স্বতঃসিদ্ধ, $x:!A \Sequent !A$।
  !!{proof}~$\delta'$ তখন অনুমান~$!A^x$।

  $n>0$ হলে $\delta$-র শেষ অনুমানবিধি
  অনুযায়ী ক্ষেত্র আলাদা করি। যেমন, সেটি
  $\Intro{\lif}$ হলে $\delta$-র শেষাংশ
  \[
  \bottomAlignProof
  \AxiomC{}
  \RightLabel{$\delta_1$}
  \Deduce$x:!B, \Gamma \fCenter !C$
  \RightLabel{\Intro{\lif}}
  \UnaryInf$\Gamma \fCenter !B \lif !C$
  \DisplayProof
  \quad\text{ বা }\quad
  \bottomAlignProof
  \AxiomC{}
  \RightLabel{$\delta_1$}
  \Deduce$\Gamma \fCenter !C$
  \RightLabel{\Intro{\lif}}
  \UnaryInf$\Gamma \fCenter !B \lif !C$
  \DisplayProof
  \]
  আরোহের অনুমানে \Log{N1c} (\Log{N1i})-তে
  $!C$-এর !!a{proof} আছে; সেটিকে $\delta_1'$ বলি।
  প্রথম ক্ষেত্রে $!B^x$ হলো~$\delta_1'$-র
  খোলা অনুমান; দ্বিতীয় ক্ষেত্রে নয়।
  % BN-SRC-558 TARGET-CORRECTION-BEGIN
  \par\smallskip\noindent\textnormal{\small\textbf{উৎস-সংশোধন (BN-SRC-558):} উৎসের বিপরীত আরোহে প্রথম ক্ষেত্রে $!B^x$-কে অন্তিম $\delta'$-র খোলা অনুমান বলা হয়েছে; →-ভূমিকায় সেটি অবমুক্ত হয়, তাই খোলা থাকে উপপ্রমাণ $\delta_1'$-তে।}
  % BN-SRC-558 TARGET-CORRECTION-END
  তাই !!{proof}~$\delta'$ যথাক্রমে নেওয়া যায়
  \[
  \bottomAlignProof
  \AxiomC{$\Discharge{!B}{x}$}
  \RightLabel{$\delta_1'$}
  \DeduceC{$!C$}
  \DischargeRule{\Intro{\lif}}{x}
  \UnaryInfC{$!B \lif !C$}
  \DisplayProof
  \quad\text{ বা }\quad
  \bottomAlignProof
  \AxiomC{}
  \RightLabel{$\delta_1'$}
  \DeduceC{$!C$}
  \RightLabel{\Intro{\lif}}
  \UnaryInfC{$!B \lif !C$}
  \DisplayProof
  \]
  % BN-SRC-557 TARGET-CORRECTION-BEGIN
  \par\smallskip\noindent\textnormal{\small\textbf{উৎস-সংশোধন (BN-SRC-557):} বিপরীত আরোহে N1-এর উপপ্রমাণটি $\delta_1'$; উৎসের শেষ দুই বৃক্ষে ভুলে N2 উপপ্রমাণ $\delta_1$-র লেবেল আছে। এখানে দুই লেবেলেই $\delta_1'$।}
  % BN-SRC-557 TARGET-CORRECTION-END
% BN-SRC-875: two displayed alternative connectors localized to Bengali.
\end{proof}

\end{document}
